\section{Filtros de suavizado}

    \subsection{Subtitulo}

        Desarrollar el ejercicio por subsecciones queda prolijo \dots

    \subsection{Preguntas y análisis}

        \begin{itemize}[leftmargin=*]
            \item \textbf{¿Qué filtro resulta más efectivo para ruido gaussiano? ¿Y para ruido sal y pimienta? Explicar por qué, relacionando el tipo de operación (lineal vs. no lineal) con la naturaleza estadística del ruido.}
            
                Desarrollar la respuesta acá \dots
            
            \item \textbf{¿Cómo afecta el tamaño del kernel al nivel de suavizado? ¿Existe un tamaño óptimo? ¿Qué se pierde al aumentar excesivamente el tamaño del kernel?}
            
                Desarrollar la respuesta acá \dots
            
            \item \textbf{Comparar visualmente y cuantitativamente (PSNR, SSIM) las imágenes filtradas. ¿Existe correlación entre los valores de estas métricas y la percepción visual de calidad? ¿En qué casos podrían divergir?}
            
                Desarrollar la respuesta acá \dots
            
            \item \textbf{El filtro de mediana, al ser no lineal, ¿presenta alguna ventaja adicional respecto a la preservación de bordes en comparación con el filtro de promedio?}
            
                Desarrollar la respuesta acá \dots
            
            \item \textbf{¿Por qué el filtro gaussiano, a diferencia del filtro de promedio simple, produce un suavizado más ``natural'' sin introducir artefactos de bloque?} 
            
                Desarrollar la respuesta acá \dots
        \end{itemize}
